% Final report. The numbers, tables and figures come from report_numbers.tex, which
% white_advantage writes when run with --report (see README.md in the repo root).
% Build from this folder:  pdflatex -jobname=finalReport report.tex  (twice)
\documentclass[11pt,a4paper]{article}
\usepackage[a4paper,margin=2cm,top=1.8cm,bottom=2cm]{geometry}
\usepackage[T1]{fontenc}
\usepackage[sc]{mathpazo}
\linespread{1.05}
\usepackage{courier}
\usepackage[scaled=0.9]{helvet}
\usepackage{microtype}
\usepackage{xcolor}
\usepackage{graphicx}
\usepackage{booktabs}
\usepackage{tabularx}
\usepackage{enumitem}
\usepackage{amsmath}
\usepackage{caption}
\usepackage[hidelinks]{hyperref}
\usepackage{titlesec}
\usepackage{placeins}
\usepackage{pgfplots}
\pgfplotsset{compat=1.16}

\titleformat{\section}{\large\sffamily\bfseries}{\thesection.}{0.5em}{}
\titleformat{\subsection}{\normalsize\sffamily\bfseries}{\thesubsection}{0.5em}{}
\titlespacing*{\section}{0pt}{12pt}{4pt}
\titlespacing*{\subsection}{0pt}{8pt}{3pt}
\setlist{itemsep=2pt,topsep=3pt,leftmargin=1.3em}
\setlength{\parindent}{0pt}
\setlength{\parskip}{4pt}
\captionsetup{font=small,labelfont=bf}

% red placeholder for anything not filled in yet
\newcommand{\TBD}[1]{\textcolor{red}{\textbf{[#1]}}}

% placeholder: white_advantage --report report_numbers.tex replaces this file with the real numbers
\newcommand{\MonthList}{\TBD{MonthList}}
\newcommand{\GamesRead}{\TBD{GamesRead}}
\newcommand{\GamesUsed}{\TBD{GamesUsed}}
\newcommand{\PctUsed}{\TBD{PctUsed}}
\newcommand{\SkipNoResult}{\TBD{SkipNoResult}}
\newcommand{\SkipAbandoned}{\TBD{SkipAbandoned}}
\newcommand{\SkipNoRating}{\TBD{SkipNoRating}}
\newcommand{\SkipGap}{\TBD{SkipGap}}
\newcommand{\FileRows}{\multicolumn{4}{c}{\TBD{filled in by white\_advantage}} \\}
\newcommand{\HAll}{\TBD{HAll}}
\newcommand{\HAllCI}{\TBD{HAllCI}}
\newcommand{\ScoreAll}{\TBD{ScoreAll}}
\newcommand{\ExpectedAll}{\TBD{ExpectedAll}}
\newcommand{\WhiteAll}{\TBD{WhiteAll}}
\newcommand{\DrawAll}{\TBD{DrawAll}}
\newcommand{\BlackAll}{\TBD{BlackAll}}
\newcommand{\FigResults}{\fbox{\parbox[c][4cm][c]{0.9\linewidth}{\centering\TBD{figure FigResults, drawn by white\_advantage}}}}
\newcommand{\OpeningRows}{\multicolumn{7}{c}{\TBD{filled in by white\_advantage}} \\}
\newcommand{\OpeningDefRows}{\multicolumn{4}{c}{\TBD{filled in by white\_advantage}} \\}
\newcommand{\OpeningGames}{\TBD{OpeningGames}}
\newcommand{\OpeningShare}{\TBD{OpeningShare}}
\newcommand{\GroupRows}{\multicolumn{7}{c}{\TBD{filled in by white\_advantage}} \\}
\newcommand{\CompareRows}{\multicolumn{6}{c}{\TBD{filled in by white\_advantage}} \\}
\newcommand{\CompareDiff}{\TBD{CompareDiff}}
\newcommand{\CompareZ}{\TBD{CompareZ}}
\newcommand{\CompareP}{\TBD{CompareP}}
\newcommand{\FigOpenings}{\fbox{\parbox[c][4cm][c]{0.9\linewidth}{\centering\TBD{figure FigOpenings, drawn by white\_advantage}}}}
\newcommand{\FigGroupH}{\fbox{\parbox[c][4cm][c]{0.9\linewidth}{\centering\TBD{figure FigGroupH, drawn by white\_advantage}}}}
\newcommand{\FigGroupDraws}{\fbox{\parbox[c][4cm][c]{0.9\linewidth}{\centering\TBD{figure FigGroupDraws, drawn by white\_advantage}}}}
\newcommand{\SpeedRows}{\multicolumn{7}{c}{\TBD{filled in by white\_advantage}} \\}
\newcommand{\SpeedStepRows}{\multicolumn{6}{c}{\TBD{filled in by white\_advantage}} \\}
\newcommand{\FigSpeed}{\fbox{\parbox[c][4cm][c]{0.9\linewidth}{\centering\TBD{figure FigSpeed, drawn by white\_advantage}}}}
\newcommand{\FigSpeedBars}{\fbox{\parbox[c][4cm][c]{0.9\linewidth}{\centering\TBD{figure FigSpeedBars, drawn by white\_advantage}}}}
\newcommand{\FigSpeedDraws}{\fbox{\parbox[c][4cm][c]{0.9\linewidth}{\centering\TBD{figure FigSpeedDraws, drawn by white\_advantage}}}}


\begin{document}

\begin{center}
{\LARGE\bfseries The Perks of Being a First Mover}\\[4pt]
{\large Does White's Advantage in Chess Change with Rating, Opening and Time Control?}\\[6pt]
CS215 Open Project \quad$\cdot$\quad Final Report \quad$\cdot$\quad Group 02\\[2pt]
{\small Sanidhya (25b0944), Daksh Maahor (25b0974), Sahil (25b0991), VadaPavJai (25b1030)}\\[2pt]
{\small Mentor: Namish Bansal}\\[2pt]
{\small\url{https://github.com/cs215-cse-iitb/dai-aug-2026-openproject-group-2}}
\end{center}

%=====================================================================
% SECTION 1: META INFORMATION
%=====================================================================

\section*{Section 1: Meta Information}

\renewcommand{\arraystretch}{1.35}
\begin{tabularx}{\linewidth}{@{}p{4.6cm}X@{}}
\toprule
\textbf{1. Open Project Title} & The Perks of Being a First Mover: Does White's Advantage in Chess Change with Rating, Opening and Time Control? \\
\textbf{2. Group ID} & ``02'' \\
\textbf{3. Mentor Name} & Namish Bansal \\
\textbf{4. Group Members and Roll Numbers} &
  Sanidhya (25b0944); Daksh Maahor (25b0974); Sahil (25b0991); VadaPavJai (25b1030) \\
\textbf{5. Honor Code and Discussions} &
  None \\
\textbf{6. LLM Usage Links} &
  \TBD{Readable links to all chat sessions / prompts used. If an AI agent was used, export its report as PDF, host it and link it here.} \\
\textbf{7. Video Presentation} & \TBD{link to video} \\
\textbf{8. GitHub Repository} & \url{https://github.com/cs215-cse-iitb/dai-aug-2026-openproject-group-2} \\
\textbf{9. GitHub Handles and Initial Commit} & \\
\hspace{1em}Sanidhya & https://github.com/SS3642 \quad earliest commit: October 21, 2026 \\
\hspace{1em}Daksh Maahor & https://github.com/Daksh-Maahor \quad earliest commit: November 3, 2026 \\
\hspace{1em}Sahil & https://github.com/Sahil-Deo-IITB \quad earliest commit: November 3, 2026 \\
\hspace{1em}VadaPavJai & https://github.com/nauseated-potato \quad earliest commit: October 18, 2026 \\
\textbf{10. Contribution Breakdown} &
  Unequal \\
\bottomrule
\end{tabularx}

\vspace{8pt}
\textbf{Work split (for reference).}
\begin{tabularx}{\linewidth}{@{}lX@{}}
Sanidhya & Statistical modelling: fitting $h$, confidence intervals and the $z$-tests. \\
Daksh Maahor & Data cleaning and opening classification; win/draw/loss rates. \\
Sahil & Data engineering: the C++ streaming parser and the Colab pipeline. \\
VadaPavJai & Visualisation: figures, the Plotly dashboard and the report. \\
\end{tabularx}

\newpage

%=====================================================================
% SECTION 2: TECHNICAL CONTENT
%=====================================================================

\section*{Section 2: Technical Content}

\subsection*{Core Problem Statement}
We measure the first-move advantage on Lichess. Instead of raw win rates, we fit White's
advantage as a number of \emph{rating points}, $h$, after correcting each game for the rating
gap between the players. Under the Elo model a player $d$ points above the opponent scores
$E(d)=1/(1+10^{-d/400})$; we give White an extra $h$ points,
$E_i(h)=1/(1+10^{-(R^{\mathrm W}_i+h-R^{\mathrm B}_i)/400})$, and choose the $h$ that makes
expected scores equal actual scores ($S_i=1,\tfrac12,0$). This is the maximum-likelihood
estimate (Newton's method; standard error from the Fisher information). Overall,
$h = \HAll \pm \HAllCI$ rating points. We test three hypotheses:
\begin{itemize}
\item \textbf{H1 (Rating).} $h$ peaks at intermediate ratings and shrinks at 2200+.
\item \textbf{H2 (Openings).} Four aggressive 1.e4 gambits (King's, Evans, Danish, Smith-Morra)
  carry a larger $h$ than four drawish openings (Petrov, Berlin, QGD, French Exchange), and the
  drawish ones have more draws at every rating.
\item \textbf{H3 (Time control).} $h$ grows from Bullet to Classical.
\end{itemize}

\subsection*{Data Source(s)}
Lichess Open Database (\url{https://database.lichess.org}, CC0), standard rated games for
\MonthList: one zstd-compressed PGN per month, tens of GB each, openly downloadable. We use
seven of ~17 header tags: \texttt{Result}, \texttt{WhiteElo}, \texttt{BlackElo},
\texttt{Event}, \texttt{TimeControl}, \texttt{Termination}, \texttt{Opening}. We read
\GamesRead{} games and keep \GamesUsed{} (\PctUsed).

\begin{center}\footnotesize
\begin{tabular}{@{}lrrr@{}}
\toprule Month & Games read & Games kept & Kept \\ \midrule
\FileRows \\ \midrule
Total & \GamesRead & \GamesUsed & \PctUsed \\ \bottomrule
\end{tabular}
\end{center}

\subsection*{Tools and Software Packages}
A C++ program streams the zstd files, reads only header lines, and counts games by opening,
speed, 100-point rating bracket and rating gap (two months on two threads, run on Google Colab,
streaming directly from Lichess). It writes \texttt{counts.csv} (a few MB); all tables, tests,
figures, this report (\LaTeX{} with pgfplots) and an interactive Plotly dashboard
(\texttt{output/final/dashboard.html}) are rebuilt from it in seconds.
\TBD{list exact Python/R packages and versions used for the analysis, if any}

% India Focus: struck out in the official guidelines -> intentionally omitted.

\subsection*{Challenges and Results Achieved}
\textbf{Cleaning and design.} We drop games with no result (\SkipNoResult), abandoned games
(\SkipAbandoned), missing ratings (\SkipNoRating) and rating gaps above 100 (\SkipGap). Games go
into 100-point brackets by average rating and into a speed from \texttt{Event}. Openings are
matched by Lichess name prefix (excluding gambit side-lines inside the drawish openings) and
cover \OpeningGames{} games (\OpeningShare{}). With millions of games every difference has a
tiny $p$-value, so we report effect sizes with 95\% CIs and treat a few rating points as small.
\TBD{technical hurdles: e.g.\ file size/streaming, opening-name matching, memory, runtime}

\textbf{Overall.} White wins \WhiteAll, draws \DrawAll{}, loses \BlackAll; average score
\ScoreAll{} versus \ExpectedAll{} expected from ratings alone, i.e.\ $h=\HAll\pm\HAllCI$.

\textbf{H1: rating.} \TBD{is H1 supported? Describe the shape of $h$ against rating}

\begin{figure}[!htbp]
\centering
\begin{minipage}[t]{0.49\linewidth}\centering
\resizebox{\linewidth}{!}{\FigSpeed}
\captionof{figure}{$h$ per 100-point bracket, by speed (95\% CIs; $<$1,000 games omitted).}
\end{minipage}\hfill
\begin{minipage}[t]{0.49\linewidth}\centering
\resizebox{\linewidth}{!}{\FigGroupH}
\captionof{figure}{$h$ by rating: aggressive vs drawish openings.}\label{fig:grouph}
\end{minipage}
\end{figure}

\textbf{H2: openings.} Each opening is fitted separately; groups are compared with a $z$-test,
$\mathrm{SE}(\hat h_A-\hat h_D)=\sqrt{\mathrm{SE}_A^2+\mathrm{SE}_D^2}$. Aggressive minus
drawish: \CompareDiff{} points ($z=\CompareZ$, $p=$ \CompareP). \TBD{is H2 supported?}

\begin{table}[!htbp]
\centering\footnotesize
\caption{White's advantage by opening (all speeds).}
\begin{tabular}{@{}lrrrrrr@{}}
\toprule
Opening & Games & Avg rating & White & Draw & Black & $h$ (pts) \\
\midrule
\OpeningRows \\ \midrule \GroupRows \\
\bottomrule
\end{tabular}
\end{table}

\textbf{H3: time control.} Each speed is its own set of games with its own $h$; each speed is
compared with the next faster one by the same $z$-test. \TBD{is H3 supported?}

\begin{table}[!htbp]
\centering\footnotesize
\caption{White's advantage by speed (all openings).}
\begin{tabular}{@{}lrrrrrr@{}}
\toprule
Speed & Games & Avg rating & White & Draw & Black & $h$ (pts) \\
\midrule
\SpeedRows \\
\bottomrule
\end{tabular}
\end{table}

\begin{figure}[!htbp]
\centering
\begin{minipage}[t]{0.49\linewidth}\centering
\resizebox{\linewidth}{!}{\FigOpenings}
\captionof{figure}{$h$ per opening with 95\% CIs.}
\end{minipage}\hfill
\begin{minipage}[t]{0.49\linewidth}\centering
\resizebox{\linewidth}{!}{\FigSpeedBars}
\captionof{figure}{$h$ per speed with 95\% CIs.}
\end{minipage}
\end{figure}

\textbf{Limitations.} Opening choice is not random (association, not causation); the two opening
groups have different players (see Fig.~\ref{fig:grouph}); ratings are noisy and Lichess is not
FIDE; games are not independent (active players contribute thousands), so true uncertainty is
somewhat larger than our CIs; one site and two months only.

\textbf{Conclusion.} \TBD{answer the three questions in a few sentences}

% Optional extra figures/tables from report_numbers.tex if space allows:
% \FigResults  \FigGroupDraws  \FigSpeedDraws  \CompareRows  \SpeedStepRows  \OpeningDefRows

\newpage

\section*{Section 3: Frequently Answered Questions (FAQ)}

\textbf{Q1. How many games did you analyse and where did they come from?}\\
\GamesUsed{} standard rated Lichess games from \MonthList{} (\GamesRead{} read, \PctUsed{} kept),
from the CC0 Lichess Open Database.

\textbf{Q2. What is $h$ and what is its overall value?}\\
$h$ is the number of Elo rating points that, added to White's rating, makes the Elo-expected
scores match the observed scores. Overall $h=\HAll\pm\HAllCI$.

\textbf{Q3. Which games did you exclude and why?}\\
Games with no result, abandoned games, missing ratings and rating gaps above 100 points: very
lopsided games say little about the first move.

\textbf{Q4. Why rating points instead of raw win rates?}\\
A raw win rate mixes the first move with the rating gap and the draw rate. Rating points correct
each game for the gap, so groups with different rating mixes are comparable, and draws count
as half a point.

\textbf{Q5. Why not fit a single slope of $h$ against rating?}\\
The shape is not close to a straight line, so a slope would hide it; we fit $h$ per
100-point bracket and read the shape from the plot.

\textbf{Q6. With millions of games, isn't everything ``significant''?}\\
Yes, so we report effect sizes with 95\% CIs and treat differences of a few rating points as
small whatever the $p$-value.

\textbf{Q7. Which openings were compared?}\\
Aggressive: King's, Evans, Danish and Smith-Morra Gambits. Drawish: Petrov, Berlin, Queen's
Gambit Declined, French Exchange (gambit side-lines excluded).

\textbf{Q8. Does the data support the hypotheses?}\\
\TBD{H1, H2, H3 each: supported / partly / not supported, with the key number}

\textbf{Q9. Does the opening \emph{cause} the difference in $h$?}\\
Not necessarily. Players choose openings, and styles, preparation and ratings differ between
groups, so our results show association only.

\textbf{Q10. What would you do with more time?}\\
\TBD{e.g.\ more months, cluster-robust CIs by player, adjusting for player ratings, FIDE/OTB data}

\newpage

\section*{Section 4: Reflection Essay (Optional)}

\TBD{Write your reflection here. Suggested prompts: key takeaways and skills learned;
how the project went beyond textbook material; why real-world problems matter in a curriculum;
relevance to placements and careers; navigating an open-ended, unguided problem; effort
compared with standard homework; personal anecdotes, challenges and final thoughts.}

\vspace{4cm}
\hrule\vspace{1cm}\hrule\vspace{1cm}\hrule


\newpage
\section*{Summary}
We measure the first-move advantage on Lichess using \GamesUsed{} standard rated games from
\MonthList. Instead of comparing raw win rates, we fit White's advantage as a number of
\emph{rating points}, $h$, after correcting every game for the rating gap between the players.
Overall White is worth $h = \HAll \pm \HAllCI$ rating points. We then ask three questions:
how $h$ changes with playing strength, with the time control (Bullet to Classical), and with
the opening, where we compare four openings with an aggressive reputation (all 1.e4 gambits)
against four with a drawish reputation.
\TBD{one or two sentences with the main findings, written once the real results are in}

\section{Question and hypotheses}
Chess is a game of perfect information in which White moves first. Database statistics have
long shown that White scores a little better than Black, but a raw win rate mixes three things:
the first move itself, the rating gap between the players, and how often games are drawn.
We separate them and test:
\begin{enumerate}[label=\textbf{H\arabic*}]
\item \textbf{Rating (from our proposal).} White's advantage is not constant across skill levels.
  It peaks at intermediate ratings, where players can exploit the initiative but defend poorly,
  and shrinks at 2200+, where Black knows how to neutralise it.
\item \textbf{Openings (new for the final stage).} Aggressive openings, where White gives up
  material for the initiative, carry a larger White advantage than drawish openings, and drawish
  openings have far more draws at every rating.
\item \textbf{Time control.} White's advantage depends on how much time the players have. In
  slower games the first-move initiative is used more accurately and fewer games are decided by
  the clock, so we expect $h$ to grow from Bullet to Classical.
\end{enumerate}

\FloatBarrier
\section{Data}
\textbf{Source.} The Lichess Open Database (\url{https://database.lichess.org}), standard rated
games for \MonthList, released under the Creative Commons CC0 licence. Each month is one
zstd-compressed PGN file of tens of GB. Every game has about 17 header tags. We use seven:
\texttt{Result}, \texttt{WhiteElo}, \texttt{BlackElo}, \texttt{Event} (gives the speed),
\texttt{TimeControl}, \texttt{Termination} and \texttt{Opening}.

\begin{center}
\small
\begin{tabular}{@{}lrrr@{}}
\toprule
Month & Games read & Games kept & Kept \\
\midrule
\FileRows
\midrule
Total & \GamesRead & \GamesUsed & \PctUsed \\
\bottomrule
\end{tabular}
\end{center}

\textbf{Cleaning.} We drop games with no result (\SkipNoResult), abandoned games
(\SkipAbandoned), games with a missing rating (\SkipNoRating), and games where the two
ratings differ by more than 100 points (\SkipGap). The rating gap is corrected for inside the
model anyway, but very lopsided games say little about the first move. Each game is put in a
100-point rating bracket $[a, a{+}100)$ by the \emph{average} of the two ratings, and in a speed
(UltraBullet to Correspondence) from the \texttt{Event} tag.

\textbf{Openings.} Lichess names the opening of every game, for example
\emph{King's Gambit Accepted: Bishop's Gambit}. We match those names as shown in
Table~\ref{tab:defs}. Inside the drawish openings we leave out their gambit side-lines, such as
the Stafford Gambit in the Petrov, so that the drawish group really is quiet. The eight openings
cover \OpeningGames{} games (\OpeningShare{} of all kept games). Everything else is ``other''.

\begin{table}[!htbp]
\centering\small
\caption{Openings and the Lichess opening names they cover. A game counts if its name starts
with one of the listed names and contains none of the excluded ones.}
\label{tab:defs}
\begin{tabularx}{\linewidth}{@{}llXp{4.4cm}@{}}
\toprule
Opening & Group & Lichess name starts with & Excluding \\
\midrule
\OpeningDefRows
\bottomrule
\end{tabularx}
\end{table}

\textbf{Why these eight.} The aggressive openings are White gambits: White gives up a pawn
or more for fast development and attack (King's, Evans, Danish and Smith-Morra Gambits). The
drawish openings have a reputation for symmetrical, simplified positions: the Petrov and the
Berlin (Black's choice), the Queen's Gambit Declined (Black's choice), and the French Exchange
(White's choice to simplify).

\FloatBarrier
\section{Method}
The full derivations, with worked examples, are in the separate document \emph{Theory.pdf}.

\FloatBarrier
\subsection{White's advantage in rating points}
Under the Elo model, a player rated $d$ points above the opponent is expected to score
$E(d) = 1/(1 + 10^{-d/400})$. We give White an extra $h$ points:
\[
E_i(h) = \frac{1}{1 + 10^{-(R^{\mathrm W}_i + h - R^{\mathrm B}_i)/400}},
\]
and choose the $h$ that makes the expected scores add up to the actual scores
($S_i = 1, \tfrac12, 0$ for a White win, draw or loss). This is the maximum-likelihood
estimate. We find it with Newton's method, and its standard error comes from the curvature of
the log-likelihood (the Fisher information). Reporting $h$ in rating points makes groups
comparable: a 1540 who beats a 1570 with White counts as stronger evidence than a win between two
1570s. Draws count as half a point. Draw and win rates are reported separately.

\FloatBarrier
\subsection{$h$ by rating (H1)}
We fit $h$ separately in every 100-point rating bracket (by the average of the two ratings), for
each speed, and read the shape from the chart of $h$ against rating. We don't fit a single line
through the brackets, because the shape is not close to a straight line, and a slope would hide
it.

\FloatBarrier
\subsection{Aggressive vs drawish (H2)}
Each opening is treated as its own set of games: we fit $h$ by maximum likelihood using only
the games played in that opening, which gives one $h$ (with its CI) per opening. For the
comparison we also fit one $h$ for all four aggressive openings together and one for the four
drawish ones,
and test the difference with a $z$-test: the two groups are different games, so the variances
add, $\mathrm{SE}(\hat h_A - \hat h_D) = \sqrt{\mathrm{SE}_A^2 + \mathrm{SE}_D^2}$. We do this for
all speeds together and for each of Bullet, Blitz and Rapid.

\FloatBarrier
\subsection{Time controls (H3)}
Lichess sorts every game into a speed (UltraBullet, Bullet, Blitz, Rapid, Classical,
Correspondence) by its time control. We treat each speed exactly like an opening: its games are
a separate set, with their own maximum-likelihood $h$ and win/draw/loss shares. We then compare
each speed with the next faster one using the same $z$-test for a difference.

\FloatBarrier
\subsection{Big data and significance}
With \GamesUsed{} games, almost any difference has a tiny $p$-value. So we always report the
size of an effect with its 95\% confidence interval, and treat effects of a few rating points
as small whatever their $p$-value.

\FloatBarrier
\subsection{Tools}
A C++ program streams the compressed files (zstd), reads only the header lines, and counts
games by opening, speed, rating bracket and rating gap. The two months are read on two threads.
We ran it on Google Colab, streaming the files straight from Lichess so nothing large was
downloaded. It saves those counts (\texttt{counts.csv}, a few MB), and every table, test, figure,
this report and an interactive Plotly dashboard (\texttt{output/final/dashboard.html}) are rebuilt
from them in seconds.

\FloatBarrier
\section{Results}

\FloatBarrier
\subsection{Overall}
Over all kept games White wins \WhiteAll, draws \DrawAll{} and loses \BlackAll. White's average
score is \ScoreAll{} against \ExpectedAll{} predicted by the ratings alone, which is
$h = \HAll \pm \HAllCI$ rating points.

\begin{figure}[!htbp]
\centering
\FigResults
\caption{Share of games won by White, drawn and won by Black, by rating (all speeds).}
\end{figure}

\FloatBarrier
\subsection{Does the advantage peak at intermediate ratings? (H1)}
\TBD{interpretation: is H1 supported?}

\begin{figure}[!htbp]
\centering
\FigSpeed
\caption{$h$ in each 100-point rating bracket, by speed, with 95\% CIs. Brackets with fewer than
1,000 games are left out.}
\end{figure}

\FloatBarrier
\subsection{Openings (H2)}
\begin{table}[!htbp]
\centering\small
\caption{White's advantage by opening (all speeds). Avg rating = average rating of the games.}
\begin{tabular}{@{}lrrrrrr@{}}
\toprule
Opening & Games & Avg rating & White & Draw & Black & $h$ (pts) \\
\midrule
\OpeningRows
\midrule
\GroupRows
\bottomrule
\end{tabular}
\end{table}

\begin{figure}[!htbp]
\centering
\FigOpenings
\caption{$h$ for each opening (all speeds) with 95\% CIs.}
\end{figure}

\begin{figure}[!htbp]
\centering
\FigGroupH
\caption{$h$ by rating for the aggressive and the drawish openings, with 95\% CIs.}\label{fig:grouph}
\end{figure}

\begin{figure}[!htbp]
\centering
\FigGroupDraws
\caption{Draw rate by rating for the aggressive and the drawish openings.}
\end{figure}

\begin{table}[!htbp]
\centering\small
\caption{Aggressive vs drawish openings: $h$ for each group, and the difference with its
$z$-test (rating points, 95\% CI), for all speeds and for each of the main speeds.}
\begin{tabular}{@{}lrrrrr@{}}
\toprule
Speed & Aggressive & Drawish & Difference & $z$ & $p$ \\
\midrule
\CompareRows
\bottomrule
\end{tabular}
\end{table}

Over all speeds, $h$ for the aggressive openings minus $h$ for the drawish openings is
\CompareDiff{} rating points ($z = \CompareZ$, $p = $ \CompareP).
\TBD{interpretation: is H2 supported, and what does it say about White's advantage?}


\FloatBarrier
\subsection{Time controls (H3)}
\begin{table}[!htbp]
\centering\small
\caption{White's advantage by speed (all openings). White / Draw / Black = share of games won
by White, drawn, won by Black. $h$ with 95\% CI.}
\begin{tabular}{@{}lrrrrrr@{}}
\toprule
Speed & Games & Avg rating & White & Draw & Black & $h$ (pts) \\
\midrule
\SpeedRows
\bottomrule
\end{tabular}
\end{table}

\begin{figure}[!htbp]
\centering
\FigSpeedBars
\caption{$h$ for each speed (all openings) with 95\% CIs.}
\end{figure}

\begin{table}[!htbp]
\centering\small
\caption{Each speed against the next faster one: $h$ for both, and the difference (slower minus
faster) with its $z$-test (rating points, 95\% CI).}
\begin{tabular}{@{}lrrrrr@{}}
\toprule
Speeds & Slower & Faster & Difference & $z$ & $p$ \\
\midrule
\SpeedStepRows
\bottomrule
\end{tabular}
\end{table}

\begin{figure}[!htbp]
\centering
\FigSpeedDraws
\caption{Draw rate by rating for each speed.}
\end{figure}

\TBD{interpretation: is H3 supported?}

\FloatBarrier
\section{Discussion}
\TBD{what the results mean for White's advantage, written once the real results are in}

\textbf{Limitations.}
\begin{itemize}
\item \emph{Choosing an opening is not random.} Players who play the King's Gambit or the
  Berlin differ in style and preparation, and the opening is chosen by one player (White for
  the gambits and the French Exchange, Black for the Petrov, Berlin and QGD). Our $h$
  corrects for the rating gap in each game, but the results show association, not that the
  opening causes the difference.
\item \emph{The groups have different players.} The comparison of the two opening groups doesn't
  adjust for their players' ratings. The ``Avg rating'' column and Figure~\ref{fig:grouph}
  (which compares the groups bracket by bracket) show how much that matters.
\item \emph{Ratings are noisy}, especially for new accounts, which blurs the correction for the
  rating gap. Lichess ratings also differ from FIDE ratings, so ``2200 on Lichess'' is not a
  master.
\item \emph{Games are not independent.} Active players contribute thousands of games, so the
  real uncertainty is somewhat larger than our confidence intervals.
\item \emph{One site, two months.} The results describe Lichess players in early 2024.
\end{itemize}

\FloatBarrier
\section{Conclusion}
\TBD{answer to the three questions in a few sentences}

\section*{Contributions}
\begin{tabularx}{\linewidth}{@{}lX@{}}
Sanidhya & Statistical modelling: fitting $h$, confidence intervals and the $z$-tests. \\
Daksh Maahor & Data cleaning and opening classification; win/draw/loss rates. \\
Sahil & Data engineering: the C++ streaming parser and the Colab pipeline. \\
VadaPavJai & Visualisation: figures, the Plotly dashboard and the report. \\
\end{tabularx}

\section*{References}
{\small
Lichess Open Database, \url{https://database.lichess.org}.
Lichess opening names, \url{https://github.com/lichess-org/chess-openings}.
A.\,E.~Elo, \emph{The Rating of Chessplayers, Past and Present}, Arco, 1978.
M.\,E.~Glickman, ``Example of the Glicko-2 system'', 2012.
}

\end{document}
